\section{Joint graphs and support cuts}\label{sec:setting}

Let $D\subset\R^n$ be nonempty and compact, and let
$F=(f_1,\dots,f_k):D\to\R^k$ be continuous. The \emph{joint graph} of $F$
on $D$ and its convex hull are
\begin{equation}\label{eq:joint-hull}
G_D(F)=\{(x,F(x)):x\in D\},\qquad H_D(F)=\conv G_D(F).
\end{equation}
In a model, the coordinates $w_j$ of $(x,w)\in H_D(F)$ stand for the
expressions $f_j(x)$ of the same variables $x$; $D$ is the set on which
these expressions must be relaxed, typically a box intersected with
linear rows. For a direction $(a,b)\in\R^n\times\R^k$ the (lower)
\emph{support value} is
\begin{equation}\label{eq:support}
h_D(a,b)=\min_{x\in D}\{a\T x+b\T F(x)\}.
\end{equation}

\begin{proposition}[Support description]\label{prop:support}
For every $(a,b)$ and every $\beta\le h_D(a,b)$, the inequality
$a\T x+b\T w\ge\beta$ is valid for $H_D(F)$. Moreover $H_D(F)$ is the
intersection of the half-spaces $a\T x+b\T w\ge h_D(a,b)$ over all
$(a,b)$.
\end{proposition}

\begin{proof}
The inequality holds on $G_D(F)$ by definition and is preserved by
convex combinations. $G_D(F)$ is compact, so $H_D(F)$ is compact and
convex, and a point outside it can be strictly separated by a linear
functional; written as a lower bound, the separating inequality is one
of the stated half-spaces.
\end{proof}

We call $a\T x+b\T w\ge\beta$ a \emph{support cut}. Two facts make
support cuts attractive. First, they are as strong as possible for the
block: no valid inequality in these coordinates is missing from the
family. Second, they combine information that separate relaxations
discard. If each $f_j$ is relaxed on its own, the result is the
intersection of the hulls of the graphs of $(x,f_j)$, which can be much
larger than $H_D(F)$ when the $f_j$ share variables
(Example~\ref{ex:two-rows}). For $n=k=1$, the family recovers the convex
and concave envelopes of a single function.

Three tasks remain, and they organize the rest of the paper: computing or
bounding $h_D(a,b)$ (Sections~\ref{sec:certification}
and~\ref{sec:quadratic}), finding a direction $(a,b)$ that separates a
given point (Section~\ref{sec:separation}), and using the cuts in a
solver without changing the model (Sections~\ref{sec:original}
and~\ref{sec:implementation}). If integer variables occur in $D$, we
relax them to their bounds; support cuts remain valid for the mixed-integer
graph, but they need not describe its convex hull.

\subsection{Related work}\label{sec:related}

\paragraph{Factorable and simultaneous relaxations.}
Global solvers relax nonconvex models by decomposing expressions into
elementary operations and relaxing each with an auxiliary variable
\citep{McCormick1976,SmithPantelides1999,TawarmalaniSahinidis2004,BelottiEtAl2009,BestuzhevaEtAl2025}.
The resulting relaxations ignore the dependence among expressions that
share variables, and they ignore linear constraints that link those
variables \citep{ZhuHeTawarmalani2026}. Simultaneous convexification
addresses the first point. \citet{Tawarmalani2010} develops inclusion
certificates for the convex hull of several functions of the same
variables; \citet{Ballerstein2013} characterizes the hull of a vector of
functions by linear combinations, and \citet{LiersEtAl2021} separate over
it in a gas network application. \citet{HeTawarmalani2021,HeTawarmalani2022,HeTawarmalani2024}
build composite relaxations with vectors of outer functions, and
\citet{ZhuHeTawarmalani2026} treat simultaneous factorable graphs with
coupling constraints. For multilinear and bilinear terms, the gap
between termwise and joint relaxations is well studied
\citep{Rikun1997,LuedtkeNamazifarLinderoth2012,BolandEtAl2017,BaoKhajaviradSahinidisTawarmalani2015},
and solvers group quadratic terms before relaxing them
\citep{MisenerFloudas2012}. Reduced-space relaxations avoid auxiliary
variables altogether \citep{TsoukalasMitsos2014,BongartzMitsos2017,NajmanBongartzMitsos2021}.
Our support cuts are an instance of simultaneous convexification; what we
add concerns how the cuts are certified, how they enter the original
model, and which support problems can be solved exactly.

\paragraph{Lagrangian and aggregation cuts.}
Cuts in the original variables obtained by minimizing a weighted sum of
constraint functions over a block are Lagrangian cuts. They appear in
block decomposition methods for MINLP \citep{Nowak2005,KaruppiahGrossmann2008,WuMutsNowakHendrix2025},
in range reduction \citep{TawarmalaniSahinidis2004,GleixnerEtAl2017}
and in rigorous filtering \citep{DomesNeumaier2016}. The convex hull
relaxation they describe, convexification in the joint space of
variables and constraint values, is the primal form of the Lagrangian
dual \citep{Falk1969,Geoffrion1974,FeltenmarkKiwiel2000,LemarechalRenaud2001},
and the set version of this statement for Lagrangian cuts is given by
\citet{ChenLuedtke2022}. Aggregations of quadratic constraints that
describe convex hulls under additional hypotheses are a different object
\citep{DeyMunozSerrano2022,BlekhermanDeySun2024}.

\paragraph{Exact low-dimensional and structured quadratic hulls.}
The convex hull of $\{(x,xx\T)\}$ over a simplex, a box in dimension two,
and small triangulated polytopes is described by semidefinite and RLT
constraints \citep{AnstreicherBurer2010}; for boxes in higher dimension
it is not \citep{BurerLetchford2009,Anstreicher2012}. Minimizing a
quadratic over a polytope of fixed dimension by enumerating faces is
classical \citep[Section~2.9]{Murty1997}, quadratic programming is in NP
\citep{Vavasis1990}, and box-constrained quadratic programming on forests
is strongly polynomial \citep{DelPiaKhajavirad2026}. Explicit hulls for
star-shaped sign patterns are given by \citet{DeyKhajavirad2025} and
\citet{Khajavirad2026}. Gluing local relaxations along a sparsity pattern
is exact when the local measures share full marginals
\citep{Vorobev1962,Lasserre2006}, and finite truncations need extra
conditions \citep{FantuzziFuentes2025}; sparse relaxations on paths can be
strictly weaker than dense ones \citep{NieDemmel2009,NieQuTangZhang2026}.
Section~\ref{sec:composition} places our results in this context.

\paragraph{Rigorous relaxations and safe cuts.}
Safe bounds and cuts for linear and mixed-integer programs correct
floating-point results with variable bounds and directed rounding
\citep{NeumaierShcherbina2004,CookEtAl2009SafeCuts,EiflerGleixner2024};
exact MIP solvers and VIPR certificates check complete solves against
the original instance \citep{CookKochSteffyWolter2013,CheungGleixnerSteffy2017,EiflerGleixner2023},
and SCIP~10 offers such an exact mode for MILP \citep{SCIP10}. For nonlinear
problems, rigorous linear relaxations and safe linear estimators are
established \citep{BorradaileVanHentenryck2005,Kearfott2011,DomesNeumaier2012,NininMessineHansen2015},
as are verified Bernstein bounds \citep{GarloffJanssonSmith2003JCAM,GarloffSmith2008,MunozNarkawicz2013}
and ball arithmetic \citep{Johansson2017arb}. That rounding in the problem
formulation itself must be controlled was pointed out by
\citet[Section~20]{Neumaier2004} and \citet{SchichlNeumaier2005}.
\citet{LiersEtAl2021} mention safe rounding of coefficients without
details. We use these techniques; to our knowledge, no published MINLP
implementation certifies the exported rows of joint support cuts against
the source model and the row stored by the solver.

\paragraph{Separation and oracles.}
Weak separation and weak optimization are polynomially equivalent for
convex bodies given by oracles \citep{GrotschelLovaszSchrijver1981,GrotschelLovaszSchrijver1988}.
Gilbert's nearest-point method gives an oracle procedure that returns
either a nearby point of the set or a separating hyperplane
\citep{Gilbert1966,BrierleyNavascuesVertesi2016}, and the local-cut and
Fenchel-cut schemes separate by column generation over an optimization
oracle, in exact arithmetic when needed \citep{Boyd1994,ChvatalCookEspinoza2013}.
Oracle-based cuts for polynomial optimization are studied by
\citet{BienstockChenMunoz2020}. Section~\ref{sec:separation} relates our
finite procedure to these results.

\paragraph{Native SCIP.}
SCIP's nonlinear constraint handler detects quadratic, bilinear,
convex, concave and other structures, uses projections of linear rows for
bilinear terms \citep{MullerSerranoGleixner2020}, separates RLT and
semidefinite minor cuts, and presolves variables that occur only
concavely to binaries \citep{BestuzhevaEtAl2025,SCIP8report}. Several valid
cut families for nonconvex problems remain disabled by default because
their overall effect is negative \citep{BestuzhevaEtAl2025,SCIP10}. A new
separator therefore has to be evaluated against this baseline, not
against termwise McCormick relaxations.

